\chapter{Introduction}\label{chap:introduction}

\section{What is \hologo{LaTeX}?}\label{sec:what-is-latex}

If you've ever had to deal with a large Word document, you know it can be a pain to do simple tasks such as updating all its references, replacing or moving an image, or just navigating through the document in general.
\Gls{wysiwyg} programs such as Word provide an instant visual preview of the document,
but they tend to reach their limits once the structure of a document becomes more important as its size grows.
Especially when wanting to change the formatting of a document later on,
it tends to rely more on the user checking that e.g.\ all the image captions look the same,
or the new heading formatting has been applied to all the headings and not just the one you happened to have selected.

This is exactly where \hologo{LaTeX} comes into play.
Instead of providing direct visual feedback to the user, it relies on a syntax that describes the structure of a document.
The user then writes plain-text files with some special annotations such as \code|\section{...}|.
These annotations mark what each part of the document is, e.g.\ a heading, a list or an equation.
\hologo{LaTeX} then uses these annotations to treat the text marked by them differently and apply a specific formatting.
This formatting is defined once for the whole document.
If you change it, every heading, caption or list in the document follows.

\section{Strengths and weaknesses}\label{sec:strengths-weaknesses}

Neither approach is better in every case.
Modern word processors also have styles, automatic captions and an equation editor,
and used with discipline, they can handle a lot.
The difference is that a word processor lets you format text directly,
while \hologo{LaTeX} makes you describe the structure first.

\begin{description}
    \item[\Gls{wysiwyg}]
        You see the result while you type, and you can start without learning anything new.
        It is also easy to work with other people, e.g.\ with comments and tracked changes.
        On the other hand, a consistent layout depends on you.
        Large documents with many images can become slow and hard to manage.
        The files are binary, so version control such as Git cannot show what changed.
    \item[\hologo{LaTeX}]
        The formatting stays consistent, because it is defined in one place.
        Numbers of sections, figures, tables and equations are set automatically,
        and so are references to them, the table of contents and the bibliography.
        \hologo{LaTeX} is very good at typesetting mathematics.
        The source files are plain text, so any editor can open them and Git can show every change.
        On the other hand, you must compile the document to see the result.
        It takes time to learn, and the error messages are hard to read at first.
        Placing one element at an exact position can take some effort,
        and everyone who edits the source must know some \hologo{LaTeX} too.
\end{description}

\autoref{fig:effort} sketches what this means in practice.
For a short text, a word processor needs less effort.
Every figure, table and equation adds more effort in a word processor than in \hologo{LaTeX},
so for a report or a thesis, \hologo{LaTeX} soon pays off.

\begin{figure}[htbp]
    \centering
    % =====================================================================
%  Effort vs. complexity of a document, for the slides and the handout.
%  A sketch, not measured data.
%  Needs: pgfplots, hologo
%  Colours: `figaccent' (LaTeX) and `figalt' (word processor), the deck's
%  colours unless defined before \input-ing this file.
% =====================================================================
\providecolor{figaccent}{HTML}{8C195F}
\providecolor{figalt}{HTML}{746BD2}
\begin{tikzpicture}
    \begin{axis}[
        smooth, samples=100,
        axis lines=left,
        axis line style={thick},
        width=.85\textwidth,
        height=4.6cm,
        xlabel=Complexity of the document,
        ylabel=Effort,
        x label style={at={(axis description cs:1,.1)}, anchor=east},
        ytick=\empty,
        xtick={2, 4, 6, 11, 13, 16, 20, 22},
        xticklabels={Text, Figures, Tables, Maths, Plots, {Technical\\ Report}, Thesis, Book},
        xticklabel style={align=center},
        x tick label style={rotate=70, anchor=east, font=\footnotesize},
        legend style={legend pos=north west},
        no markers, domain=0:25,
    ]
        \addplot[densely dashed, very thick, figalt, domain=0:14]{exp(x/3 - 2.5)};
        \addlegendentry{Word processor}
        \addplot[very thick, figaccent]{3/(1 + exp(-2*(x-2.5))) + x/50*exp((x-4)/12)};
        \addlegendentry{\textrm{\hologo{LaTeX}}}
    \end{axis}
\end{tikzpicture}

    \caption{Effort for documents of growing complexity. A sketch, not a measurement.}
    \label{fig:effort}
\end{figure}

\section{From source file to PDF}\label{sec:source-to-pdf}

A \hologo{LaTeX} document starts as a plain-text source file with the extension \texttt{.tex}.
A program called the engine, e.g.\ \hologo{pdfLaTeX} or \hologo{LuaLaTeX}, reads this file and writes the PDF.
This step is called compiling (\autoref{fig:pipeline}).
You never edit the PDF itself, only the source file.

\begin{figure}[htbp]
    \centering
    % =====================================================================
%  From source file to PDF, for the slides and the handout.
%  Needs: tikz, hologo
%  Colours: `figaccent' and `figfill', the deck's colours unless
%  defined before \input-ing this file.
% =====================================================================
\providecolor{figaccent}{HTML}{8C195F}
\providecolor{figfill}{HTML}{F1F0FB}
\begin{tikzpicture}[
    stage/.style={
        draw=figaccent, line width=1pt, rounded corners=1.3mm,
        fill=figfill,
        minimum width=3.2cm, minimum height=1.6cm, align=center,
    },
    flow/.style={->, line width=1.2pt, figaccent},
]
    \node[stage] (src) at (0,0)   {Source file\\[1mm]\texttt{report.tex}};
    \node[stage] (tex) at (4.6,0) {\hologo{LaTeX}};
    \node[stage] (pdf) at (9.2,0) {PDF\\[1mm]\texttt{report.pdf}};
    \node[above, text=figaccent] at (tex.north) {compile};
    \draw[flow] (src) -- (tex);
    \draw[flow] (tex) -- (pdf);
    \draw[flow] (pdf.south) to[out=-120, in=-60]
        node[below, text=figaccent] {change the source, compile again} (src.south);
\end{tikzpicture}

    \caption{From source file to PDF. After every change to the source file, you compile again.}
    \label{fig:pipeline}
\end{figure}

Because the formatting is not part of the text, one small change can affect the whole document.
The two documents in \autoref{fig:class-comparison} differ in a single word, the document class:
\texttt{article} at the top, \texttt{scrartcl} at the bottom.
The heading font, the margins and the text size change, but not a single word of the text.
\autoref{sec:document-classes} explains the document class in more detail.

\begin{figure}[htbp]
    \centering
    % trim = left bottom right top, in bp: the same part of both pages,
    % from the left edge of scrartcl's text to the right edge of article's.
    \slidecode[firstline=1, lastline=1, numbers=none]{00-class-article.tex}
    \sliderender[trim=80 677 108 70, clip, width=\dimexpr\linewidth+2\exoverhang\relax]{00-class-article.pdf}
    \slidecode[firstline=1, lastline=1, numbers=none]{00-class-scrartcl.tex}
    \sliderender[trim=80 677 108 70, clip, width=\dimexpr\linewidth+2\exoverhang\relax]{00-class-scrartcl.pdf}
    \caption{The same document with two document classes. Both crops show the same part of the page.}
    \label{fig:class-comparison}
\end{figure}

\section{A short history}\label{sec:history}

\hologo{LaTeX} is not a new idea.
\autoref{fig:tex-timeline} shows how it developed.
Donald Knuth wrote \hologo{TeX}, the program \hologo{LaTeX} is based on, in 1978 \parencite{knuth1984}.
In 1984, Leslie Lamport released \hologo{LaTeX} as a set of commands on top of \hologo{TeX} that are easier to use \parencite{lamport1994}.
The version used today, \hologo{LaTeXe}, was released in 1994 and is still maintained by the \hologo{LaTeX} Project,
which publishes a new release twice a year.

\begin{figure}[htbp]
    \centering
    % One year is narrower than on the slides, so the timeline fits the text width.
    \newcommand{\timelineunit}{0.32cm}
    % =====================================================================
%  Timeline of TeX, LaTeX and the engines, for the slides and the handout.
%  Needs: tikz, hologo
%  Above the line: TeX and the macro packages written for it.
%  Below the line: newer programs (engines) that run them.
%
%  \timelineunit: width of one year. Default 0.4cm (slides);
%  define it before \input-ing this file to change it, e.g. in the handout:
%      \newcommand{\timelineunit}{0.36cm}
%  Colour: `figaccent', the deck's purple unless defined before.
% =====================================================================
\providecommand{\timelineunit}{0.4cm}
\providecolor{figaccent}{HTML}{8C195F}
\begin{tikzpicture}[
    x=\timelineunit,
    tl year/.style={font=\scriptsize\bfseries, text=figaccent},
    tl name/.style={font=\small, align=center},
    tl who/.style={font=\scriptsize, text=black!60, align=center},
]
    % 1975 at x=0, one unit per year.
    % \lvl: 1 or 2 above the line, -1 or -2 below; 2 sits further out,
    % for entries that would otherwise touch their neighbour.
    \draw[line width=1pt, figaccent, ->] (1,0) -- (35,0);
    \foreach \y/\lbl/\person/\lvl in {
            1978/\hologo{TeX}/Donald Knuth/1,
            1984/\hologo{LaTeX}/Leslie Lamport/1,
            1994/\hologo{LaTeXe}/\hologo{LaTeX}\\Project/1,
            1996/\hologo{ConTeXt}/Hans Hagen/2,
            1996/\hologo{pdfTeX}/H\`an Th\'{\^e} Th\`anh/-1,
            2004/\hologo{XeTeX}/Jonathan Kew/-1,
            2007/\hologo{LuaTeX}/\hologo{LuaTeX} team/-2} {
        \pgfmathsetmacro{\x}{\y-1975}
        \pgfmathsetmacro{\out}{(abs(\lvl)-1)*1.25}
        \draw[figaccent, line width=1pt] (\x,-.15) -- (\x,.15);
        % the year always sits under the line; same year twice is fine
        \node[tl year] at (\x,-.4) {\y};
        \ifnum\lvl>0
            \draw[black!30] (\x,.15) -- (\x,.3+\out);
            \node[tl who, anchor=south] (tl-who) at (\x,.3+\out) {\person};
            \node[tl name, anchor=south, inner ysep=1pt] at (tl-who.north) {\lbl};
        \else
            \draw[black!30] (\x,-.65) -- (\x,-.8-\out);
            \node[tl name, anchor=north, inner ysep=1pt] (tl-name) at (\x,-.8-\out) {\lbl};
            \node[tl who, anchor=north] at (tl-name.south) {\person};
        \fi
    }
\end{tikzpicture}

    \caption{\hologo{TeX} and \hologo{LaTeX} above the line, newer engines below.}
    \label{fig:tex-timeline}
\end{figure}

Newer engines were added over time.
\hologo{pdfTeX} (around 1996) writes PDF files directly,
while \hologo{XeTeX} (2004) and \hologo{LuaTeX} (2007) can use the fonts installed on your computer.
\hologo{ConTeXt} (around 1996) is a separate system based on \hologo{TeX}, but it is not a part of \hologo{LaTeX}.

In this workshop and in this handout, \enquote{compiling} always means \hologo{pdfLaTeX}.
It is the default on Overleaf and in most editors,
and everything shown here works with it.
